\chapter{Grundlagen}
\label{sec:grundlagen}

Dieses Kapitel vermittelt die für das Verständnis des Prototyps notwendigen Grundlagen. Nach einer kurzen Einordnung des Weigh-in-Motion-Konzepts liegt der Schwerpunkt auf den konkret verwendeten Werkzeugen: dem ADXL335-Beschleunigungssensor, dem zugehörigen SDP-B-Controller-Board, der Auswertesoftware sowie den beiden Hilfsprogrammen Sandboxie und AutoHotkey, die für die Durchführung der Messungen benötigt wurden.

\section{Weigh-in-Motion: Kurzüberblick}
\label{sec:grundlagen-wim}

Weigh-in-Motion (WIM) bezeichnet ein Verfahren zur Bestimmung von Achslasten und Gesamtgewichten von Fahrzeugen im fließenden Verkehr, ohne dass der Verkehrsfluss unterbrochen werden muss \cite{QFree2020}. Klassische, fahrbahnintegrierte WIM-Systeme erreichen dabei Genauigkeiten zwischen $\pm 5\,\%$ und $\pm 15\,\%$ nach COST-323, während vibrationsbasierte Ansätze wie der von Bajwa auf den angrenzenden Untergrund als Messmedium zurückgreifen \cite{QFree2020, Bajwa:EECS-2013-210}. Das vorliegende Projekt setzt an diesem vibrationsbasierten Ansatz an, ohne selbst eine vollständige WIM-Anlage nachzubilden.

\section{ADXL335-Beschleunigungssensor}
\label{sec:grundlagen-adxl335}

Der ADXL335 ist ein 3-Achsen-Beschleunigungssensor mit einem Messbereich von $\pm 3g$, der an seinen drei Ausgängen jeweils eine zur Beschleunigung proportionale analoge Spannung liefert \cite{ADXL335_Datasheet}. Er bildet die zentrale Messkomponente des im Rahmen dieses Projekts entwickelten Aufbaus zur Erfassung von Bodenvibrationen.

% Platzhalter: Foto des ADXL335-Sensors bzw. des Sensoraufbaus am Fahrbahnrand
\begin{figure}[htbp]
\centering
\includegraphics[width=0.4\textwidth]{fig/EVAL-ADXL354BZANGLE-web}
\caption{ADXL335-Beschleunigungssensor}
\label{fig:adxl335-foto}
\end{figure}

\section{SDP-B Controller Board}
\label{sec:grundlagen-sdpb}

Das SDP-B Controller Board (EVAL-SDP-CB1Z) dient als Verbindungsglied zwischen dem ADXL335-Sensor und dem PC. Es basiert auf einem Analog Devices ADSP-BF527 Blackfin-Prozessor und stellt über zwei 120-Pin-Steckverbinder die Schnittstelle zur Sensor-Platine bereit, während die Kommunikation mit dem PC über eine USB-2.0-Verbindung erfolgt \cite{ADI_UG277}. Der Sensor wird über die 120-Pin-Verbindung mechanisch und elektrisch an das Board angeschlossen; das Board selbst wird anschließend per USB mit dem Laptop verbunden, auf dem die Auswertesoftware läuft \cite{ADI_UG277}.

% Platzhalter: Foto des SDP-B Boards bzw. der Steckverbindung Sensor-Board
\begin{figure}[htbp]
\centering
% \includegraphics[width=0.5\textwidth]{fig/sdp-b-board}
\caption{SDP-B Controller Board mit angeschlossenem Sensor}
\label{fig:sdpb-foto}
\end{figure}

\section{Auswertesoftware und Bedienoberfläche}
\label{sec:grundlagen-software}

% Platzhalter: Beschreibung der EVAL-ADXL35x-SDP-Software (GUI),
% Erklärung Messstart und einstellbarer Parameter (z. B. Sample-Rate, Messbereich, Dauer)
Die Auswertung der Sensordaten erfolgt über die zugehörige Bedienoberfläche der EVAL-ADXL35x-SDP-Plattform \cite{ADI_UG2190}. Über diese Oberfläche wird die Messung gestartet und Parameter wie Abtastrate, Messbereich und Aufzeichnungsdauer eingestellt (Platzhalter: genaue Parameterwerte gemäß später erhobener Konfiguration ergänzen).

\begin{figure}[htbp]
\centering
% \includegraphics[width=0.8\textwidth]{bilder/eval_gui_screenshot.png}
\caption{Bedienoberfläche der Auswertesoftware}
\label{fig:gui-screenshot}
\end{figure}

\section{Messaufbau und Workflow}
\label{sec:grundlagen-workflow}

Der grundlegende Messaufbau ergibt sich aus dem Zusammenspiel der zuvor beschriebenen Komponenten: Der ADXL335-Sensor wird an das SDP-B-Board angeschlossen (Abschnitt~\ref{sec:grundlagen-sdpb}), welches wiederum über USB mit dem Laptop verbunden wird \cite{ADI_UG277}. Auf dem Laptop läuft die Auswertesoftware (Abschnitt~\ref{sec:grundlagen-software}), über die zunächst die Messparameter festgelegt werden, bevor die eigentliche Messung gestartet wird \cite{ADI_UG2190}. Nach Abschluss der Messung können sowohl die Rohmesswerte als auch eine grafische Darstellung des zeitlichen Signalverlaufs aus der Software exportiert werden. Dieser Ablauf -- Anschluss, Konfiguration, Messung, Export -- bildet den grundlegenden Workflow, der in den nachfolgenden Kapiteln bei jedem Testzyklus wiederkehrt.

\section{Sandboxie}
\label{sec:grundlagen-sandboxie}

Sandboxie ist eine Isolationssoftware für Windows, die es ermöglicht, Anwendungen in voneinander getrennten, virtuellen Umgebungen (Sandboxes) parallel auszuführen \cite{Sandboxie2026}. Im Rahmen dieses Projekts war der Einsatz von Sandboxie notwendig, um für den später beschriebenen Dual-Sensor-Aufbau zwei Instanzen der Auswertesoftware gleichzeitig und unabhängig voneinander betreiben zu können; die konkrete Umsetzung wird in Abschnitt~\ref{sec:realisierung-dualsensor} beschrieben.

\section{AutoHotkey}
\label{sec:grundlagen-autohotkey}

AutoHotkey ist eine frei verfügbare Scripting-Sprache für Windows, mit der sich wiederkehrende Bedienabläufe automatisieren und per Tastenkürzel auslösen lassen \cite{AutoHotkey2026}. Durch den Einsatz von AutoHotkey-Skripten konnten wiederkehrende Bedienschritte der Messsoftware automatisiert werden, wodurch sich die Effizienz der Messdurchläufe im Projektverlauf deutlich erhöhte.
% Grobgerüst Realisierung -- stichpunktartig, noch nicht ausformuliert

\chapter{Realisierung}
\label{sec:realisierung}

% Aufbau: iterative Testzyklen statt starrer Themen-Unterkapitel
% Zyklusmuster pro Abschnitt:
% 1. Nachdenken / Meeting -> Erkenntnisse aus vorherigem Zyklus
% 2. Ideen / Probleme auflisten
% 3. Umsetzung / Test
% 4. Ergebnis -> Übergang zum nächsten Zyklus

\section{Zyklus 1: [Kurztitel, z. B. Grundaufbau Einzelsensor]}
\label{sec:realisierung-zyklus1}
% - Ausgangspunkt: nur Sensor + Software (siehe Ausgangssituation, sec:einleitung-ausgangssituation)
% - Anwendung des Grundworkflows aus sec:grundlagen-workflow
% - Erste Testmessung, grundlegende Funktionsprüfung
% - Erkannte Probleme / offene Fragen
% - Daraus abgeleitete nächste Schritte

\section{Zyklus 2: [Kurztitel]}
\label{sec:realisierung-zyklus2}
% - Platzhalter für zweiten Iterationsschritt

\section{Zyklus n: Dual-Sensor-Aufbau mit Sandboxie}
\label{sec:realisierung-dualsensor}
% - Problem: zwei Sensor-Instanzen gleichzeitig betreiben
% - Lösung: Sandboxie (siehe sec:grundlagen-sandboxie)
% - Testdurchführung Dual-Sensor
% - Ergebnis / offene Punkte

\section{Effizienzsteigerung durch AutoHotkey-Automatisierung}
\label{sec:realisierung-autohotkey}
% - Ausgangsproblem: manuelle Bedienschritte pro Messung zu langsam/fehleranfällig
% - Umsetzung: AutoHotkey-Skripte (siehe sec:grundlagen-autohotkey)
% - Effekt auf Messdurchsatz / Datenqualität

\chapter{Ergebnisse}
\label{sec:ergebnisse}
\enquote{Neuigkeiten} Messergebnisse
